\chapter*{Hướng Dẫn Sử Dụng Sách}
\addcontentsline{toc}{chapter}{Hướng Dẫn Sử Dụng Sách}
\markboth{Hướng Dẫn Sử Dụng Sách}{Hướng Dẫn Sử Dụng Sách}

\begin{truyen}{Quyển Sách Này Dành Cho Ai?}{Who Is This Book For?}
\chuhoa{Q}{uyển sách này dành cho học sinh.}
\textit{(This book is for students.)}

Nó cũng dành cho phụ huynh.
\textit{(It is also for parents.)}

Nó dành cho giáo viên.
\textit{(It is for teachers.)}

Và thành thật hơn, nó dành cho bất kỳ người lớn nào từng nói câu ``điện thoại hại con lắm'' trong khi chính mình cũng không buông nổi.
\textit{(And more honestly, it is for any adult who has said ``phones are harming children'' while being unable to put theirs down either.)}
\end{truyen}

\begin{baihoc}
Đọc quyển này tốt nhất là mỗi lần một chương.
\textit{(The best way to read this book is one chapter at a time.)}

Đọc xong, đừng hỏi ngay người khác có giống thế không.
\textit{(After reading, do not immediately ask whether others are like this.)}

Hãy hỏi hôm nay mình đã cúi xuống màn hình nhiều đến mức nào.
\textit{(Ask instead how many times today you yourself bent down toward the screen.)}
\end{baihoc}

\begin{tuhocnhanh}
1. Em muốn dùng quyển này để nhắc ai nhất? \textit{(Whom do you most want this book to remind?)}

2. Em có đủ công bằng để để mình vào cùng cái gương với người khác không? \textit{(Are you fair enough to place yourself before the same mirror as others?)}

3. Nếu phải bắt đầu sửa đúng một thói quen dùng điện thoại, em sẽ sửa điều gì trước? \textit{(If you had to start fixing one phone habit, what would you change first?)}
\end{tuhocnhanh}

\ngancach